\section{Hypersonic Overview\label{ch2:hypersonicOverview}}
        The first hypersonic flight took place during the early stages of the
        Cold War, in $1949$ when the German V2 rocket developed during the
        Second World War in Nazi Germany and brought to the USA (possible
        during operation Paperclip) reached hypersonic speeds as it entered the
        atmosphere \cite{Anderson2022}. In $1961$, the Soviet Union cosmonaut
        Yuri Garin became the fist person to safely reach space and orbit
        Earth. During his journey he became the first human in history to
        travel at hypersonic speeds \cite{Anderson2022}.\\

        As the velocity of the flow increases, the compressibility effects can
        not be ignored. The flow field can be characterized in four regimes (subsonic,
        transonic, supersonic and hypersonic).
        \begin{itemize}
            \item \textbf{Subsonic:} A flow characterized by a Mach number less
                than $1$ at every point. In this regime, compressibility effects can generally be neglected. The flow patterns in a subsonic flow are typically smooth and continuous, with no shock wave formation. Due to the absence of shock waves, these flows are often considered isentropic.
            \item \textbf{Transonic:} A flow characterized by a Mach number
                ranging between $0.8$ to $1.2$. In this regime, regions of subsonic and
                supersonic coexist; creating complex flow patterns and leading
                to the formation of localized shock waves. Compressibility
                effects cannot be ignored, leading to variations in fluid
                density.
            \item \textbf{Supersonic:} A flow characterized by a Mach number
                greater than $1$. In this regime compressibility effects
                cannot be ignored, leading to substantial variations in fluid
                density. Shock waves are present in this regime which produces
                discontinuous changes in pressure, temperature and density. 
            \item \textbf{Hypersonic:} A flow characterized by a Mach number
                greater than $5$. In this regime strong shock waves are
                present, leading to abrupt and significant changes in pressure,
                temperature, and density. Due to the high velocities, the
                kinetic energy in the flow is converted to thermal energy
                resulting in high temperatures. These high temperatures lead to
                real gas effects and a thicker boundary layer.
        \end{itemize}

        \FigRef{fig:ch2ShockMachRegions} show these regions, as the Mach number
        increases, shock waves start to form and compressibility effects become
        more important.

        \pic{0.85}{chapter2/shockMachRegions}{Mach number
        regions}{fig:ch2ShockMachRegions}

        Speeds in a hypersonic flow are greater than Mach five, at these speeds
        the kinetic energy of the gas molecules is converted into thermal
        energy, leading to very high temperatures. At these speeds the chemical
        composition of the gas becomes important. The high temperature can
        cause dissociation (breaking apart of molecules), and ionization
        (formation of charged particles) of the air molecules. At temperatures above $2000\Units{K}$ $O_2$ starts to dissociate, and
        it is completely dissociated above $4000\Units{K}$. $N_2$ begins to
        dissociate above $4000\Units{K}$ and it is completely
        dissociated above $9000\Units{K}$ \cite{Anderson2006}. \\

        At these temperatures, the Hamiltonian (kinetic and potential energy)
        of the system has to be modified because more energy modes are
        accessible (rotational, vibrational, and electronic energy).
        \FigRef{fig:ch2AdditionalDoFs} shows the additional degrees of freedom that
        emerge during a hypersonic flow.

        \pic{0.9}{chapter2/aditionalDoFs}{Additional Degrees of
        Freedom}{fig:ch2AdditionalDoFs}

        In addition, each new accessible energy mode acts at different time scales, \EqRef{timeScales} show these time scales.
        \begin{equation}\label{timeScales}
            \tau_{flow} \approx \tau_{tr} < \tau_{rot} < \tau_{vib} < \tau_{chem}
        \end{equation}

        The different time scales is the reason why flows at hypersonic
        speeds are in nonequilibrium. After air molecules travel across a
        shock, it will take certain number of collisions to convert the
        energy into thermal energy. In the case of translational energy is
        almost instantaneous, in the case of rotational is approximately $10$
        collisions and in the case of vibrational is approximately $1000$
        collisions \cite{Vincenti1975}. The nonequilibrium
        characteristics of the flow can be characterize by comparing the flow
        characteristic time with the other time scales \cite{Neitzel2017}. \\

        These additional energy modes increase the \gls{dof} of the system,
        allowing molecules to occupy more energy levels (rotational,
        vibrational, electronic). Due to this, the partition function (which
        indicates how particles in an assembly are partitioned among various
        energy levels of an atom or molecule \cite{Laurendeau2005}) has to be
        modified to consider these additional \glspl{dof}. In statistical
        mechanics, the partition function is used to calculate the heat
        capacities of the system, which are then used to calculate macro
        properties such as heat transfer and temperature, which are employed in
        fluid dynamics. \FigRef{fig:ch2HeatCapacity} shows the change in heat capacity when
        additional \glspl{dof} are accessible.

        \pic{1.0}{chapter2/heatCapacity}{Heat Capacity}{fig:ch2HeatCapacity}

        At these conditions, the calorically perfect gas approximation is no
        longer suitable because this assumption does not consider thermochemical
        effects. \FigRef{fig:caloricallyPerfectGas} shows the temperature pre
        and post shock for different altitudes. As expected at high speeds, the
        calorically perfect gas under-predicts temperature across the shock.

        \pic{0.85}{chapter2/temperatureRatio}{Temperature ratio across a shock
        wave\cite{Anderson2006}}{fig:caloricallyPerfectGas}

      Therefore, at these speeds, thermochemical effects need to be considered and accurate models that consider these effects need to be included. \\
      \todo{Add more about chemical effects affecting fluid properties}
